\chapter*{About This Edition}
\addcontentsline{toc}{chapter}{About This Edition}

The full edition of this book runs to nearly three hundred pages because its argument is a systems argument, and a systems argument has to be shown layer by layer before it can be believed: the model, the compute, the semiconductor, the watt, the data, physical production, trust, governance, finance, and the international order that all of them together imply. Many readers do not need the layers. They need the shape---what is being claimed, why, what would prove it wrong, and what it means for a government, a laboratory, or a company that is neither in Washington nor in Beijing. This condensed edition is for them.

Three rules govern the abridgement. First, nothing here is new: every number, grade, forecast, and recommendation is taken from the full edition of the same version and date, which carries the sources, the evidence grades, the derivations, and the falsification apparatus. Where the condensed text states a figure, the full edition states where it came from and how much to trust it. Second, the industrial precedents that occupy Part~I of the full edition---the solar, battery, and electric-vehicle wars, and the campaigns where the same playbook has failed---are here summarized in a single chapter, because for the purposes of the argument what matters is the playbook they establish and its limits, not the campaign histories. Third, the weight of this edition is shifted deliberately toward the overall strategy and its international consequences: what the convergence means for the West, for Japan and Europe, for the non-aligned majority of the world, and for the institution the analysis keeps arriving at. Readers who want the engineering---the memory-architecture arithmetic, the appliance cost model, the sizing framework, the security fault tree---will find pointers into the full edition throughout, marked \full{n}.

The title is deliberately provocative, but it is not a prediction of inevitability; it is a structural analysis. ``Winning'' means what it meant in solar, batteries, and electric vehicles: capturing the overwhelming majority of global production, setting the cost curve, controlling the supply chain, defining the de facto standards, and reducing erstwhile leaders to protected niches. By that definition China has already won three technology wars in which the West began with every advantage. The claim of this book is that artificial intelligence is following the same trajectory, and that the mechanism is visible only when the elements are laid side by side.

A note on the author's position is owed and is given in full in the preface to the full edition. In brief: I direct the RIKEN Center for Computational Science, whose flagship program runs on an American accelerator architecture and whose deepest partnerships are with American and European institutions. Nothing in my professional position gives me an incentive to argue that the Chinese stack is winning; a great deal of it gives me an incentive to argue the opposite. I did not begin this project expecting to arrive where it has arrived, and I have tried to hold the argument to a standard of proof appropriate to how much I would prefer it to be wrong. Where the evidence does not support the strong claim, the book---and this abridgement of it---says so.

The currency date of this edition is \datafreeze, the date on the title page. Every factual claim is current as of that date; anything learned after it belongs to the next edition. This is a living document. Corrections are welcome, and the fastest way to change its conclusions is to falsify an entry in the full edition's evidence ledger.
